% E8: the same book, with the phases replaced by a work pool.
% Every number here is a macro from paper/e8_results.tex, generated by
% experiments/e8_adaptive_book/analyze.py from the sealed run under
% runs/e8-rawapi-p16.

\subsection{E8: A Population That Does Not Wait for a Phase}
\label{sec:eval-e8}

E7's population moves through phases, and at every boundary it waits for its
slowest member. At $p=16$ on one machine that idle time was
\eSevenCoordShareSixteen\% of rank-time, and the reason is not the protocol: a
population of \eEightRanks\ agents whose slowest model takes four times the
median will wait for that model at every barrier, however cheap the barrier is.
E8 asks how much of the wait is the phase structure rather than the models, and
what a runtime has to offer before a harness author can write the alternative.

\paragraph{What changed.} The book is cut into pages rather than into $p$
segments, and every rank owns a contiguous \emph{home block} of about six of
them. After the glossary is bound, there are no phases: a rank that finishes a
page claims the next from a \emph{work pool} (Section~\ref{sec:win}), its own
block first and then the block with the most pages left; a seam between two
finished pages enters the pool the moment both exist and is claimed by whichever
rank is free; a translator that settles a new term publishes it with one atomic
union, and every later translation, anyone's, reads it. The root publishes the
settled glossary with a nonblocking broadcast and goes straight to the pool. The
only collective after the glossary is the reduction that totals the spend.

Everything else is E7's: the same book, the same model pool, the same prompts,
contracts, in-conversation repair, memoisation and epochs. The comparison is
therefore between two harnesses over one runtime, which is the comparison the
paper wants to be able to make.

\paragraph{The protocol had to grow twice, and both are general.} Writing E8
against the runtime as it stood was not possible, and what was missing was not
book-shaped. A bag of tasks over shared state gets one of four things wrong on
its own --- two workers take the same item, a dead worker keeps its item, an
item runs before its inputs exist, nobody can say when the job is finished ---
and none of the four has any application in it. They became \code{Pool\_*}
(S9.5, Section~\ref{sec:win}). A root that must wait for its receivers before
returning to its own work became \code{Ibcast} (S7.4, Section~\ref{sec:coll}).
Both are in the conformance suite on every transport and in the command binding.
What stayed with the harness is exactly what is about books: that an item is a
page or a seam, that a seam depends on two pages, that a rank prefers its home
block, and what a translation is. That line --- the runtime owns the obligations
that are the same in every program --- is the same one the rest of the design
draws, and E8 is the first experiment that tests whether it falls in a usable
place.

\paragraph{What it cost and what it bought.} The configuration was run twice.
\eEightRanks\ ranks over two machines finished the book in \eEightWall\ and
\eEightBWall\ minutes for \$\eEightCost\ and \$\eEightBCost, rendering
\eEightCoverage\% and \eEightBCoverage\% of the \eEightParagraphs\ paragraphs,
with all \eEightSeams\ seams revised each time and \eEightStolen\ and
\eEightBStolen\ pages taken from another rank's block. Pages per rank ran from
\eEightPagesMin\ to \eEightPagesMax\ against a mean of \eEightPagesMean: the pool
balanced a population whose models differ by a factor of four in speed, which a
static partition cannot do.

Model work was \eEightModelH\ and \eEightBModelH\ rank-hours against E7's
\eSevenWorkHSixteen\ at the same scale, as it must be, since it is the same book
through the same models. What the pool changed is the rest, and the two runs only
agree once the waiting is split in two. E7 spent \eSevenBlockedHSixteen\
rank-hours blocked at its barriers. E8 spent \eEightWaitWorkH\ and
\eEightBWaitWorkH\ rank-hours waiting for work \emph{while work existed} --- the
idleness the pool was built to remove, and it is gone, by an order of magnitude,
in both runs.

\paragraph{What a pool cannot remove.} The two runs' totals differ by a factor of
two, \eEightIdleShare\% of rank-time idle against \eEightBIdleShare\%, and all of
the difference is in one place. A pool cannot finish before its last item, so
once one item remains the rest of the population has nothing to do however long
that item takes. In the first run the last item took \eEightTailMin\ minutes and
cost \eEightWaitTailH\ rank-hours; in the second a single seam's model call ran
for \eEightBSlowestCallS\ seconds --- returning normally, not failing --- and the
tail cost \eEightBWaitTailH\ rank-hours, more than the whole population's model
work. The call also outlasted the rank's lease, so a peer reclaimed the item and
translated the seam a second time; both copies were correct and the later write
won, which is the pool's reclaim rule behaving exactly as specified and paying
for a straggler twice.

This is the honest limit of the result. Against a heavy-tailed executor the pool
removes the idleness that comes from structure and cannot touch the idleness that
comes from the tail, and at $p=16$ with 189 items the tail is one call. A harness
that wanted it back would have to replicate the last items speculatively, which
is a scheduling policy and not a protocol mechanism; the trace prices it either
way, which is the property that matters here.

\paragraph{What the run showed that the design did not predict.} Two machines
sharing one compare-and-swap ref are not two peers. The first E8 launch stalled
one machine for ten minutes with no failure anywhere: its ranks sat waiting for
their daemon to acknowledge a trace append, and the daemon lost every push race
to the other machine, whose translators were pushing back to back. The loser's
backoff grew with each defeat while the winner never paused, so the loser's
fetch-commit-push never fitted in a gap. Two changes followed, both in the
device: a writer that has just pushed, or that sees another machine's commits
land, yields for longer than the other's round trip; and a trace append is
acknowledged when queued rather than when pushed, so a rank's program never
waits on its own evidence. The specification now permits the second explicitly
(S13) and requires that reading the trace land it. Neither is about pools; both
were invisible until a harness stopped inserting barriers, which had been
serialising the writers by accident.
